\documentclass[a4paper,12pt]{amsbook}
\usepackage{physics}
\usepackage{amsmath,amssymb,amsthm}
\usepackage{mathtools}
\usepackage{mathpazo}
\usepackage{microtype}
\providecommand{\tb}{\textbackslash}
\newcommand{\0}{{$|0\rangle$}}
\newcommand{\1}{{$|1\rangle$}}
\usepackage{graphicx}
\usepackage{amsfonts}
\usepackage[utf8]{inputenc}
\usepackage[T1]{fontenc}
\usepackage{geometry}
\geometry{a4paper, margin=1in}
\usepackage{listings}
\usepackage{xcolor}
\pagecolor{black}
\color{white}
\usepackage{parskip}
\setlength{\parindent}{0pt}

\begin{document}
\author{Vivek Soorya Maadoori}
\title{Quantun Operations, Quantum Noise}
\maketitle
\tableofcontents

\section{Quantum Operations}

\vspace{1cm}
\textbf{\underline{Exercise 8.1}: (Unitary evolution as a quantum operation)} Pure states evolve under unitary transforms as $|\psi\rangle \rightarrow U|\psi\rangle$. Show that, equivalently, we may write $\rho \rightarrow \varepsilon (\rho) \equiv U\rho U^\dag \text{, for } \rho = |\psi\rangle\langle \psi|$.

\underline{Solution}:
\begin{enumerate}
    \item If pure states evolve as $|\psi\rangle \rightarrow U|\psi\rangle$, then their duals evolve like so $\langle \psi| \rightarrow \langle \psi| U^\dag$.

    \item Now conider how $\rho$ would evolve:

    $ \rho = |psi\rangle\langle psi| \rightarrow U |psi\rangle \langle psi | U^\dag = U \rho U^\dag$.
\end{enumerate}

$\blacksquare$

\vspace{1cm}
\textbf{\underline{Exercise 8.2}: (Measurement as a quantum operation)}

Recall from Section 2.2.3 (on page 84) that a quantum measurement with outcomes labeled by \textit{m} is described by a set of measurement operators $M_m$ such that $$\Sigma_m M^\dag_m M_m = I.$$ Let the state of the system immediately before the measurement be $\rho$. Show that for $$\varepsilon_m(\rho) \equiv M_m \rho M_m^\dag,$$ the state of the system immediately after the measurement is $$\frac{\varepsilon_m (\rho)}{tr(\varepsilon(\rho))}.$$ Also show that the probability of obtaining this measurement result is p(m) = tr$(\varepsilon_m (\rho))$.

\vspace{0.5cm}
\underline{Solution}:
First lets show the post measurement state of the density operator,
\begin{enumerate}

    \item Note that $\varepsilon_m(\rho) \equiv M_m \rho M_m^\dag = M_m |\psi\rangle\langle\psi|M_m^\dag $

    \item Post measurement state is $\frac{M_m |\psi\rangle}{\sqrt{\langle\psi| M_m^\dag M_m |\psi\rangle}}$

    \item Post measurement stat of the dual is  $\frac{\langle \psi | M_m^\dag}{\sqrt{\langle\psi| M_m M_m^\dag |\psi\rangle}}$

    \item Thus, the post measurement state of $\rho$ is $\frac{M_m |\psi\rangle\langle \psi | M_m^\dag}{\langle\psi| M_m^\dag M_m |\psi\rangle}$ \footnote{Assuming the measurement operator commutes with its adjoint. Does it?}

    \item Note that the numerator is $\varepsilon_m(\rho)$ and that the denominator is $tr(\varepsilon_m(\rho))$.
\end{enumerate}

\vspace{0.5cm}
Now, lets see the probability
\begin{enumerate}
    \item Note that the probability of measuring m after measurement is given by $p(m) = \langle\psi| M_m^\dag M_m |\psi\rangle$
    \item Note that this is $tr(\varepsilon(\rho))$
\end{enumerate}
$\blacksquare$
\end{document}
